
\title{\textbf{Mathematical Foundations, Architecture Diagrams,\\and State Analysis of Modern LLM Training Methods}}
\author{Chief Strategist J}
\date{\today}
\maketitle
\tableofcontents
\newpage

%% ============================================================
\part{Foundation Training Methods}
%% ============================================================

\section{Pretraining}

\subsection{Mathematical Foundation}
Causal language modeling models the probability of a sequence $x = (x_1, \ldots, x_T)$ as the product of autoregressive conditional probabilities:
\begin{equation}
  P(x) = \prod_{t=1}^T P(x_t \mid x_{<t})
\end{equation}
The hidden representation $h_t^{(l)}$ at layer $l$ and token position $t$ is computed using multi-head causal self-attention. The causal attention mask matrix $M \in \mathbb{R}^{T \times T}$ is defined as:
\begin{equation}
  M_{i,j} = \begin{cases} 0 & \text{if } j \le i \\ -\infty & \text{if } j > i \end{cases}
\end{equation}
This mask is added directly to the scaled query-key dot products:
\begin{equation}
  \text{Attention}(Q, K, V) = \text{softmax}\left( \frac{Q K^\top}{\sqrt{d_k}} + M \right) V
\end{equation}
where $Q = X W_Q$, $K = X W_K$, $V = X W_V$. This prevents information leakage from future tokens. At the output layer, the logits $z_t \in \mathbb{R}^{|\mathcal{V}|}$ are computed via $z_t = h_t^{(L)} W_{\text{vocab}}$. To maintain numerical stability, softmax is computed by subtracting the maximum logit:
\begin{equation}
  P(x_t = v \mid x_{<t}) = \frac{\exp(z_{t,v} - \max_k z_{t,k})}{\sum_{j \in \mathcal{V}} \exp(z_{t,j} - \max_k z_{t,k})}
\end{equation}
The causal model parameters $\theta$ are trained by minimizing the cross-entropy loss over a dataset $\mathcal{D}$:
\begin{equation}
  \mathcal{L}_{\text{PT}}(\theta) = -\frac{1}{|\mathcal{D}|}\sum_{x \in \mathcal{D}} \sum_{t=1}^T \log P_\theta(x_t \mid x_{<t})
\end{equation}
Updates are made using the AdamW optimizer. Let $g_t = \nabla_\theta \mathcal{L}_t$ be the gradient at step $t$. The update equations are:
\begin{align}
  m_t &= \beta_1 m_{t-1} + (1-\beta_1) g_t \\
  v_t &= \beta_2 v_{t-1} + (1-\beta_2) g_t^2 \\
  \hat{m}_t &= \frac{m_t}{1-\beta_1^t}, \quad \hat{v}_t = \frac{v_t}{1-\beta_2^t} \\
  \theta_t &= \theta_{t-1} - \eta \left( \frac{\hat{m}_t}{\sqrt{\hat{v}_t} + \epsilon} + \lambda \theta_{t-1} \right)
\end{align}
where $\eta$ is the learning rate, $\lambda$ is the weight decay rate, $\beta_1, \beta_2$ are the exponential decay rates, and $\epsilon$ prevents division by zero. A cosine learning rate schedule with linear warmup is employed:
\begin{equation}
  \eta_t = \eta_{\text{min}} + \frac{1}{2}(\eta_{\text{max}} - \eta_{\text{min}})\left(1 + \cos\left(\frac{t - T_{\text{warmup}}}{T_{\text{max}} - T_{\text{warmup}}}\pi\right)\right)
\end{equation}

\subsection{Architecture Flow Diagram}

\begin{figure}[h]
\centering
\begin{tikzpicture}[node distance=0.9cm and 1.4cm]
  \node[datblk] (corpus) {Raw Text\\Corpus $\mathcal{D}$};
  \node[grayblk, right=of corpus] (tok) {Tokenise\\$x_{1..T}$};
  \node[frzblk, right=of tok] (fwd) {Transformer\\Forward Pass};
  \node[lossblk, right=of fwd] (loss) {Cross-Entropy\\Loss $\mathcal{L}$};
  \node[trnblk, below=0.8cm of loss] (bwd) {Backprop +\\Adam Step};
  \node[mrgblk, left=3.6cm of bwd] (model) {Base Model\\$\theta^*$};

  \draw[arr] (corpus)--(tok);
  \draw[arr] (tok)--(fwd);
  \draw[arr] (fwd)--(loss);
  \draw[arr] (loss)--(bwd);
  \draw[arr] (bwd)--node[below,font=\tiny]{update $\theta$}(model);
  \draw[arr] (model.north) |- ($(fwd.south)+(0,-0.3)$) -- (fwd.south);
\end{tikzpicture}
\caption{Pretraining pipeline: corpus sampling, forward pass, loss computation, and weight update loop.}
\end{figure}

\subsection{State Transition Table}
\begin{table}[h]\centering
\begin{tabular}{lll}\toprule
\textbf{Property} & \textbf{Before} & \textbf{After}\\\midrule
Parameter source & Random init $\theta_0 \sim \mathcal{N}(0,\sigma^2)$ & Converged weights $\theta^*$\\
Training corpus & Not consumed & Trillions of tokens processed\\
Perplexity & Very high ($>10^3$) & Low ($<20$ on held-out data)\\
World knowledge & None & Broad language \& factual knowledge\\
requires\_grad & True (all params) & True (all params)\\\bottomrule
\end{tabular}
\caption{State properties before and after pretraining.}
\end{table}

\newpage
%% ------------------------------------------------------------
\section{Continued Pretraining (CPT)}

\subsection{Mathematical Foundation}
Causal language modeling models the probability of a sequence $x = (x_1, \ldots, x_T)$ as the product of autoregressive conditional probabilities:
\begin{equation}
  P(x) = \prod_{t=1}^T P(x_t \mid x_{<t})
\end{equation}
The hidden representation $h_t^{(l)}$ at layer $l$ and token position $t$ is computed using multi-head causal self-attention. The causal attention mask matrix $M \in \mathbb{R}^{T \times T}$ is defined as:
\begin{equation}
  M_{i,j} = \begin{cases} 0 & \text{if } j \le i \\ -\infty & \text{if } j > i \end{cases}
\end{equation}
This mask is added directly to the scaled query-key dot products:
\begin{equation}
  \text{Attention}(Q, K, V) = \text{softmax}\left( \frac{Q K^\top}{\sqrt{d_k}} + M \right) V
\end{equation}
where $Q = X W_Q$, $K = X W_K$, $V = X W_V$. This prevents information leakage from future tokens. At the output layer, the logits $z_t \in \mathbb{R}^{|\mathcal{V}|}$ are computed via $z_t = h_t^{(L)} W_{\text{vocab}}$. To maintain numerical stability, softmax is computed by subtracting the maximum logit:
\begin{equation}
  P(x_t = v \mid x_{<t}) = \frac{\exp(z_{t,v} - \max_k z_{t,k})}{\sum_{j \in \mathcal{V}} \exp(z_{t,j} - \max_k z_{t,k})}
\end{equation}
The causal model parameters $\theta$ are trained by minimizing the cross-entropy loss over a dataset $\mathcal{D}$:
\begin{equation}
  \mathcal{L}_{\text{PT}}(\theta) = -\frac{1}{|\mathcal{D}|}\sum_{x \in \mathcal{D}} \sum_{t=1}^T \log P_\theta(x_t \mid x_{<t})
\end{equation}
Updates are made using the AdamW optimizer. Let $g_t = \nabla_\theta \mathcal{L}_t$ be the gradient at step $t$. The update equations are:
\begin{align}
  m_t &= \beta_1 m_{t-1} + (1-\beta_1) g_t \\
  v_t &= \beta_2 v_{t-1} + (1-\beta_2) g_t^2 \\
  \hat{m}_t &= \frac{m_t}{1-\beta_1^t}, \quad \hat{v}_t = \frac{v_t}{1-\beta_2^t} \\
  \theta_t &= \theta_{t-1} - \eta \left( \frac{\hat{m}_t}{\sqrt{\hat{v}_t} + \epsilon} + \lambda \theta_{t-1} \right)
\end{align}
where $\eta$ is the learning rate, $\lambda$ is the weight decay rate, $\beta_1, \beta_2$ are the exponential decay rates, and $\epsilon$ prevents division by zero. A cosine learning rate schedule with linear warmup is employed:
\begin{equation}
  \eta_t = \eta_{\text{min}} + \frac{1}{2}(\eta_{\text{max}} - \eta_{\text{min}})\left(1 + \cos\left(\frac{t - T_{\text{warmup}}}{T_{\text{max}} - T_{\text{warmup}}}\pi\right)\right)
\end{equation}

\subsection{Architecture Flow Diagram}

\begin{figure}[h]
\centering
\begin{tikzpicture}[node distance=0.9cm and 1.3cm]
  \node[frzblk] (base) {Pretrained\\Base $\theta_{\text{base}}$};
  \node[datblk, above right=0.6cm and 1.4cm of base] (new) {New Corpus\\$\mathcal{D}_{\text{new}}$};
  \node[grayblk, below right=0.6cm and 1.4cm of base] (replay) {Replay Buffer\\$\mathcal{D}_{\text{replay}}$};
  \node[grayblk, right=3.2cm of base] (mix) {Batch\\Mixer $\lambda$};
  \node[lossblk, right=of mix] (loss) {LM Loss\\$\mathcal{L}_{\text{CPT}}$};
  \node[trnblk, right=of loss] (upd) {Gradient\\Update};
  \node[mrgblk, below=1.1cm of upd] (out) {Adapted\\Model $\theta'$};

  \draw[arr] (new.east) -| (mix.north);
  \draw[arr] (replay.east) -| (mix.south);
  \draw[arr] (base) -- (mix);
  \draw[arr] (mix) -- (loss);
  \draw[arr] (loss) -- (upd);
  \draw[arr] (upd) -- (out);
\end{tikzpicture}
\caption{CPT mixes new corpus and replay buffer batches to update a pretrained checkpoint.}
\end{figure}

\subsection{State Transition Table}
\begin{table}[h]\centering
\begin{tabular}{lll}\toprule
\textbf{Property} & \textbf{Before} & \textbf{After}\\\midrule
Starting checkpoint & Pretrained $\theta_{\text{base}}$ & $\theta' = \theta_{\text{base}} + \Delta\theta$\\
New domain knowledge & Absent & Incorporated\\
Original knowledge & Retained & Partially retained via replay\\
Learning rate schedule & Finished & Warm-restart applied\\
Perplexity on new domain & High & Significantly reduced\\\bottomrule
\end{tabular}
\caption{State properties before and after continued pretraining.}
\end{table}

\newpage
%% ------------------------------------------------------------
\section{Domain Adaptive Pretraining (DAPT)}

\subsection{Mathematical Foundation}
Causal language modeling models the probability of a sequence $x = (x_1, \ldots, x_T)$ as the product of autoregressive conditional probabilities:
\begin{equation}
  P(x) = \prod_{t=1}^T P(x_t \mid x_{<t})
\end{equation}
The hidden representation $h_t^{(l)}$ at layer $l$ and token position $t$ is computed using multi-head causal self-attention. The causal attention mask matrix $M \in \mathbb{R}^{T \times T}$ is defined as:
\begin{equation}
  M_{i,j} = \begin{cases} 0 & \text{if } j \le i \\ -\infty & \text{if } j > i \end{cases}
\end{equation}
This mask is added directly to the scaled query-key dot products:
\begin{equation}
  \text{Attention}(Q, K, V) = \text{softmax}\left( \frac{Q K^\top}{\sqrt{d_k}} + M \right) V
\end{equation}
where $Q = X W_Q$, $K = X W_K$, $V = X W_V$. This prevents information leakage from future tokens. At the output layer, the logits $z_t \in \mathbb{R}^{|\mathcal{V}|}$ are computed via $z_t = h_t^{(L)} W_{\text{vocab}}$. To maintain numerical stability, softmax is computed by subtracting the maximum logit:
\begin{equation}
  P(x_t = v \mid x_{<t}) = \frac{\exp(z_{t,v} - \max_k z_{t,k})}{\sum_{j \in \mathcal{V}} \exp(z_{t,j} - \max_k z_{t,k})}
\end{equation}
The causal model parameters $\theta$ are trained by minimizing the cross-entropy loss over a dataset $\mathcal{D}$:
\begin{equation}
  \mathcal{L}_{\text{PT}}(\theta) = -\frac{1}{|\mathcal{D}|}\sum_{x \in \mathcal{D}} \sum_{t=1}^T \log P_\theta(x_t \mid x_{<t})
\end{equation}
Updates are made using the AdamW optimizer. Let $g_t = \nabla_\theta \mathcal{L}_t$ be the gradient at step $t$. The update equations are:
\begin{align}
  m_t &= \beta_1 m_{t-1} + (1-\beta_1) g_t \\
  v_t &= \beta_2 v_{t-1} + (1-\beta_2) g_t^2 \\
  \hat{m}_t &= \frac{m_t}{1-\beta_1^t}, \quad \hat{v}_t = \frac{v_t}{1-\beta_2^t} \\
  \theta_t &= \theta_{t-1} - \eta \left( \frac{\hat{m}_t}{\sqrt{\hat{v}_t} + \epsilon} + \lambda \theta_{t-1} \right)
\end{align}
where $\eta$ is the learning rate, $\lambda$ is the weight decay rate, $\beta_1, \beta_2$ are the exponential decay rates, and $\epsilon$ prevents division by zero. A cosine learning rate schedule with linear warmup is employed:
\begin{equation}
  \eta_t = \eta_{\text{min}} + \frac{1}{2}(\eta_{\text{max}} - \eta_{\text{min}})\left(1 + \cos\left(\frac{t - T_{\text{warmup}}}{T_{\text{max}} - T_{\text{warmup}}}\pi\right)\right)
\end{equation}

\subsection{Architecture Flow Diagram}

\begin{figure}[h]
\centering
\begin{tikzpicture}[node distance=0.9cm and 1.2cm]
  \node[datblk] (domain) {Domain Corpus\\(Medical/Legal/Code)};
  \node[grayblk, right=of domain] (clean) {Clean \&\\Deduplicate};
  \node[frzblk, right=of clean] (base) {Base Model\\$\theta_{\text{base}}$};
  \node[lossblk, right=of base] (loss) {Domain LM\\Loss $\mathcal{L}$};
  \node[trnblk, below=0.9cm of loss] (step) {Adam Step\\low LR};
  \node[mrgblk, left=2.2cm of step] (dapt) {DAPT Model\\$\theta_{\text{domain}}$};

  \draw[arr] (domain)--(clean);
  \draw[arr] (clean)--(base);
  \draw[arr] (base)--(loss);
  \draw[arr] (loss)--(step);
  \draw[arr] (step)--node[below,font=\tiny]{converged}(dapt);
  \draw[arr,dashed,blue!60] (dapt.north) -- ++(0,0.5) -| (base.south) node[midway,above,font=\tiny]{checkpoint};
\end{tikzpicture}
\caption{DAPT pipeline: domain corpus is cleaned, fed into the base model, and trained with standard LM loss.}
\end{figure}

\subsection{State Transition Table}
\begin{table}[h]\centering
\begin{tabular}{lll}\toprule
\textbf{Property} & \textbf{Before} & \textbf{After}\\\midrule
Domain vocabulary familiarity & Low & High\\
Domain perplexity & High ($\rho \gg 1$) & Low (near 1)\\
General language ability & Strong & Marginally reduced\\
Downstream task (same domain) & Weak & Significantly stronger\\
Training data type & General web text & Domain-specific unlabelled\\\bottomrule
\end{tabular}
\caption{State properties before and after domain adaptive pretraining.}
\end{table}

\newpage
%% ------------------------------------------------------------
\section{Task Adaptive Pretraining (TAPT)}

\subsection{Mathematical Foundation}
Causal language modeling models the probability of a sequence $x = (x_1, \ldots, x_T)$ as the product of autoregressive conditional probabilities:
\begin{equation}
  P(x) = \prod_{t=1}^T P(x_t \mid x_{<t})
\end{equation}
The hidden representation $h_t^{(l)}$ at layer $l$ and token position $t$ is computed using multi-head causal self-attention. The causal attention mask matrix $M \in \mathbb{R}^{T \times T}$ is defined as:
\begin{equation}
  M_{i,j} = \begin{cases} 0 & \text{if } j \le i \\ -\infty & \text{if } j > i \end{cases}
\end{equation}
This mask is added directly to the scaled query-key dot products:
\begin{equation}
  \text{Attention}(Q, K, V) = \text{softmax}\left( \frac{Q K^\top}{\sqrt{d_k}} + M \right) V
\end{equation}
where $Q = X W_Q$, $K = X W_K$, $V = X W_V$. This prevents information leakage from future tokens. At the output layer, the logits $z_t \in \mathbb{R}^{|\mathcal{V}|}$ are computed via $z_t = h_t^{(L)} W_{\text{vocab}}$. To maintain numerical stability, softmax is computed by subtracting the maximum logit:
\begin{equation}
  P(x_t = v \mid x_{<t}) = \frac{\exp(z_{t,v} - \max_k z_{t,k})}{\sum_{j \in \mathcal{V}} \exp(z_{t,j} - \max_k z_{t,k})}
\end{equation}
The causal model parameters $\theta$ are trained by minimizing the cross-entropy loss over a dataset $\mathcal{D}$:
\begin{equation}
  \mathcal{L}_{\text{PT}}(\theta) = -\frac{1}{|\mathcal{D}|}\sum_{x \in \mathcal{D}} \sum_{t=1}^T \log P_\theta(x_t \mid x_{<t})
\end{equation}
Updates are made using the AdamW optimizer. Let $g_t = \nabla_\theta \mathcal{L}_t$ be the gradient at step $t$. The update equations are:
\begin{align}
  m_t &= \beta_1 m_{t-1} + (1-\beta_1) g_t \\
  v_t &= \beta_2 v_{t-1} + (1-\beta_2) g_t^2 \\
  \hat{m}_t &= \frac{m_t}{1-\beta_1^t}, \quad \hat{v}_t = \frac{v_t}{1-\beta_2^t} \\
  \theta_t &= \theta_{t-1} - \eta \left( \frac{\hat{m}_t}{\sqrt{\hat{v}_t} + \epsilon} + \lambda \theta_{t-1} \right)
\end{align}
where $\eta$ is the learning rate, $\lambda$ is the weight decay rate, $\beta_1, \beta_2$ are the exponential decay rates, and $\epsilon$ prevents division by zero. A cosine learning rate schedule with linear warmup is employed:
\begin{equation}
  \eta_t = \eta_{\text{min}} + \frac{1}{2}(\eta_{\text{max}} - \eta_{\text{min}})\left(1 + \cos\left(\frac{t - T_{\text{warmup}}}{T_{\text{max}} - T_{\text{warmup}}}\pi\right)\right)
\end{equation}

\subsection{Architecture Flow Diagram}

\begin{figure}[h]
\centering
\begin{tikzpicture}[node distance=0.8cm and 1.2cm]
  \node[datblk] (labeled) {Labelled Task\\Data (small)};
  \node[grayblk, above=0.7cm of labeled] (unlabeled) {Unlabelled Task-\\Adjacent Text};
  \node[frzblk, right=2.0cm of labeled] (dapt) {DAPT / Base\\Model $\theta$};
  \node[lossblk, right=of dapt] (loss) {Task LM\\Loss};
  \node[trnblk, right=of loss] (update) {Update\\$\theta \to \theta_{\text{task}}$};
  \node[mrgblk, below=0.9cm of update] (sft) {SFT\\(next step)};

  \draw[arr] (unlabeled.east) -| (dapt.north);
  \draw[arr] (labeled.east) -- node[above,font=\tiny]{strip labels} (dapt.west);
  \draw[arr] (dapt)--(loss);
  \draw[arr] (loss)--(update);
  \draw[arr] (update)--(sft);
\end{tikzpicture}
\caption{TAPT uses task-adjacent unlabelled text to bridge base model and supervised fine-tuning.}
\end{figure}

\subsection{State Transition Table}
\begin{table}[h]\centering
\begin{tabular}{lll}\toprule
\textbf{Property} & \textbf{Before} & \textbf{After}\\\midrule
Task vocabulary alignment & Partial & Strong\\
Corpus size & Millions of docs & Thousands of docs\\
Training epochs & 1--3 & 10--100 (small corpus)\\
Downstream label efficiency & Requires more labels & Requires fewer labels\\
Transfer to other tasks & Good & Reduced (narrow focus)\\\bottomrule
\end{tabular}
\caption{State properties before and after task adaptive pretraining.}
\end{table}

\newpage
%% ============================================================
\part{Supervised Fine-Tuning Methods}
%% ============================================================

\section{Supervised Fine-Tuning (SFT)}

\subsection{Mathematical Foundation}
SFT trains the model on structured prompt-response pairs $(x, y) \sim \mathcal{D}_{\text{SFT}}$ using teacher forcing. Let $s = [x ; y]$ be the concatenated token sequence of length $M + U$, where $x = (s_1, \ldots, s_M)$ is the prompt and $y = (s_{M+1}, \ldots, s_{M+U})$ is the response. The key math lies in applying a target label mask $L \in \{-100, \mathcal{V}\}^{M+U}$ defined as:
\begin{equation}
  L_t = \begin{cases} -100 & \text{if } t \le M \\ s_t & \text{if } t > M \end{cases}
\end{equation}
The cross-entropy loss function ignores targets set to $-100$, directing gradient updates exclusively to the response generation tokens:
\begin{equation}
  \mathcal{L}_{\text{SFT}}(\theta) = -\frac{1}{U} \sum_{t=M+1}^{M+U} \log P_\theta(s_t \mid s_{<t})
\end{equation}
This loss ignores the prompt distribution $P(x)$, preventing the model from degrading its output generation capability to match arbitrary prompts.
The gradients only backpropagate through response tokens:
\begin{equation}
  \nabla_\theta \mathcal{L}_{\text{SFT}} = -\frac{1}{U} \sum_{t=M+1}^{M+U} \nabla_\theta \log P_\theta(s_t \mid s_{<t})
\end{equation}

\subsection{Architecture Flow Diagram}

\begin{figure}[h]
\centering
\begin{tikzpicture}[node distance=0.8cm and 1.2cm]
  \node[datblk] (pairs) {$(x, y)$\\Curated Pairs};
  \node[grayblk, right=of pairs] (concat) {Concat\\{[x ; y]}};
  \node[frzblk, right=of concat] (fwd) {Transformer\\Forward};
  \node[lossblk, right=of fwd] (loss) {Masked CE\\Loss $\mathcal{L}_{\text{SFT}}$};
  \node[trnblk, below=0.9cm of loss] (bwd) {Backprop\\(response only)};
  \node[mrgblk, left=3.5cm of bwd] (out) {Instruction-\\Following $\theta^*$};

  \draw[arr] (pairs)--(concat);
  \draw[arr] (concat)--(fwd);
  \draw[arr] (fwd)--(loss);
  \draw[arr] (loss)--(bwd);
  \draw[arr] (bwd)--(out);

  \node[red,font=\tiny,text width=2cm,align=center] at ($(fwd.south)+(0,-0.3)$) {prompt tokens\\masked in labels};
\end{tikzpicture}
\caption{SFT pipeline with prompt-masked loss: gradients flow only through response tokens.}
\end{figure}

\subsection{State Transition Table}
\begin{table}[h]\centering
\begin{tabular}{lll}\toprule
\textbf{Property} & \textbf{Before} & \textbf{After}\\\midrule
Behaviour & Completion only & Instruction following\\
Loss region & N/A & Response tokens only\\
Label masking & None & Prompt masked to $-100$\\
Data format & Raw text & (instruction, response) pairs\\
Output quality & Random completion & Task-aligned responses\\\bottomrule
\end{tabular}
\caption{State properties before and after supervised fine-tuning.}
\end{table}

\newpage
%% ------------------------------------------------------------
\section{Instruction Tuning}

\subsection{Mathematical Foundation}
Instruction Tuning exposes the model to multi-turn conversational patterns. A prompt template $T(I, x)$ converts an instruction $I$ and context $x$ into structured sequences. Turn boundaries are defined using chat-formatting systems (like ChatML). A single sequence layout is:
\begin{equation}
  S = \langle\text{user}\rangle I \parallel x \langle\text{assistant}\rangle y \langle\text{eos}\rangle
\end{equation}
Let $\mathcal{D}_{\text{inst}}$ consist of instructions spanning hundreds of tasks. The model parameters $\theta$ are optimized on the joint multi-task dataset:
\begin{equation}
  \mathcal{L}_{\text{IT}}(\theta) = -\mathbb{E}_{(I, x, y) \sim \mathcal{D}_{\text{inst}}} \sum_{t=1}^{|y|} \log P_\theta\bigl(y_t \mid T(I, x), y_{<t}\bigr)
\end{equation}
We mask all prompt-side tokens, including conversational format tags (like $\langle\text{user}\rangle$ and instruction text), ensuring parameters are trained only on synthesizing assistant responses:
\begin{equation}
  \mathcal{L}_{\text{IT}}(\theta) = -\frac{1}{|y|} \sum_{i=1}^{|T(I,x)|+|y|} M_i \log P_\theta(S_i \mid S_{<i}), \quad M_i = \mathbb{I}(i > |T(I, x)|)
\end{equation}

\subsection{Architecture Flow Diagram}

\begin{figure}[h]
\centering
\begin{tikzpicture}[node distance=0.7cm and 1.1cm]
  \node[datblk] (t1) {Summarise\\Task $\mathcal{T}_1$};
  \node[datblk, below=0.4cm of t1] (t2) {Translate\\Task $\mathcal{T}_2$};
  \node[datblk, below=0.4cm of t2] (t3) {QA / Code\\Task $\mathcal{T}_K$};
  \node[grayblk, right=1.6cm of t2] (tmpl) {Chat Template\\Formatter};
  \node[frzblk, right=of tmpl] (model) {LLM\\$\theta$};
  \node[lossblk, right=of model] (loss) {SFT Loss\\(multi-task)};
  \node[mrgblk, below=0.8cm of loss] (out) {Instruction-\\Tuned $\theta^*$};

  \draw[arr] (t1.east) -| (tmpl.north);
  \draw[arr] (t2)--(tmpl);
  \draw[arr] (t3.east) -| (tmpl.south);
  \draw[arr] (tmpl)--(model);
  \draw[arr] (model)--(loss);
  \draw[arr] (loss)--(out);
\end{tikzpicture}
\caption{Instruction tuning over diverse task types funnelled through a unified chat template.}
\end{figure}

\subsection{State Transition Table}
\begin{table}[h]\centering
\begin{tabular}{lll}\toprule
\textbf{Property} & \textbf{Before} & \textbf{After}\\\midrule
Zero-shot generalisation & Poor & Strong across task types\\
Task diversity in training & N/A & $K$ heterogeneous tasks\\
Input format & Free-form text & Structured chat template\\
Prompt sensitivity & High & Reduced\\
Held-out task performance & Weak & Strong (FLAN, InstructGPT)\\\bottomrule
\end{tabular}
\caption{State properties before and after instruction tuning.}
\end{table}

\newpage
%% ------------------------------------------------------------
\section{Multi-Task Training}

\subsection{Mathematical Foundation}
Multi-Task Training aggregates losses from $K$ distinct tasks. A simple summation can result in tasks with large gradient scales dominating the parameter updates. To resolve this, we use the GradNorm algorithm. Let $W$ be the shared parameter weights (e.g., the last transformer block weights). We define the gradient norm for task $k$ at training step $t$ as:
\begin{equation}
  G_k(t) = \|\nabla_W (w_k(t) \mathcal{L}_k(t))\|_2
\end{equation}
Let $\bar{G}(t) = \frac{1}{K} \sum_k G_k(t)$ be the average gradient norm. The target gradient norm for task $k$ balances task difficulties:
\begin{equation}
  \hat{G}_k(t) = \bar{G}(t) \times \left( \frac{\tilde{\mathcal{L}}_k(t)}{\frac{1}{K}\sum_j \tilde{\mathcal{L}}_j(t)} \right)^\alpha
\end{equation}
where $\tilde{\mathcal{L}}_k(t) = \mathcal{L}_k(t) / \mathcal{L}_k(0)$ is the task loss ratio indicating task training progress, and $\alpha$ is a hyperparameter managing structural gradient balancing. The task weights $w_k(t)$ are updated by minimizing the L1 distance between $G_k(t)$ and $\hat{G}_k(t)$:
\begin{equation}
  \mathcal{L}_{\text{grad}}(w_1, \dots, w_K; t) = \sum_{k=1}^K \left| G_k(t) - \hat{G}_k(t) \right|
\end{equation}
The final joint parameter objective is:
\begin{equation}
  \mathcal{L}_{\text{MT}}(\theta) = \sum_{k=1}^K w_k \mathcal{L}_k(\theta)
\end{equation}

\subsection{Architecture Flow Diagram}

\begin{figure}[h]
\centering
\begin{tikzpicture}[node distance=0.6cm and 1.1cm]
  \node[datblk] (d1) {Dataset $\mathcal{D}_1$};
  \node[datblk, below=0.4cm of d1] (d2) {Dataset $\mathcal{D}_2$};
  \node[datblk, below=0.4cm of d2] (dk) {Dataset $\mathcal{D}_K$};
  \node[frzblk, right=1.8cm of d2, minimum height=3.2cm, minimum width=3.0cm] (backbone) {Shared\\Transformer\\Backbone $\theta_{\text{shared}}$};
  \node[trnblk, right=1.6cm of d1] (h1) {Head$_1$\\$\theta_1$};
  \node[trnblk, right=1.6cm of d2] (h2) {Head$_2$\\$\theta_2$};
  \node[trnblk, right=1.6cm of dk] (hk) {Head$_K$\\$\theta_K$};
  \node[lossblk, right=1.6cm of h2] (loss) {Weighted\\Loss $\mathcal{L}_{\text{MT}}$};

  \draw[arr] (d1)--(backbone.west|-h1);
  \draw[arr] (d2)--(backbone);
  \draw[arr] (dk)--(backbone.west|-hk);
  \draw[arr] (backbone.east|-h1)--(h1);
  \draw[arr] (backbone.east|-h2)--(h2);
  \draw[arr] (backbone.east|-hk)--(hk);
  \draw[arr] (h1.east) -| (loss.north);
  \draw[arr] (h2)--(loss);
  \draw[arr] (hk.east) -| (loss.south);
\end{tikzpicture}
\caption{Multi-task training: each dataset feeds task-specific heads on a shared backbone.}
\end{figure}

\subsection{State Transition Table}
\begin{table}[h]\centering
\begin{tabular}{lll}\toprule
\textbf{Property} & \textbf{Before} & \textbf{After}\\\midrule
Backbone specialisation & Single task & $K$ tasks jointly\\
Task heads & 1 & $K$ heads\\
Generalisation & Narrow & Broad\\
Data efficiency & Per-task only & Cross-task transfer\\
Gradient conflicts & N/A & Managed via $\lambda_k$\\\bottomrule
\end{tabular}
\caption{State properties before and after multi-task training.}
\end{table}

\newpage
%% ------------------------------------------------------------
\section{Chain-of-Thought SFT}

\subsection{Mathematical Foundation}
CoT SFT trains models to output intermediate reasoning steps before generating the final answer. Let $x$ be the prompt, $c = (c_1, \dots, c_M)$ be the reasoning chain tokens, and $y = (y_1, \dots, y_U)$ be the final answer tokens. The joint probability factorization is:
\begin{equation}
  P(c, y \mid x) = \prod_{i=1}^M P_\theta(c_i \mid x, c_{<i}) \prod_{j=1}^U P_\theta(y_j \mid x, c, y_{<j})
\end{equation}
The model is trained by minimizing the joint loss:
\begin{equation}
  \mathcal{L}_{\text{CoT}}(\theta) = -\beta \sum_{i=1}^M \log P_\theta(c_i \mid x, c_{<i}) - (1 - \beta) \sum_{j=1}^U \log P_\theta(y_j \mid x, c, y_{<j})
\end{equation}
During autoregressive generation, we sample tokens using Temperature-scaled softmax and Nucleus (top-$p$) filtering. The logits $z_t$ are modified as:
\begin{equation}
  \tilde{z}_{t,i} = \frac{z_{t,i}}{T}
\end{equation}
where $T > 0$ is the temperature. The top-$p$ vocabulary subset $\mathcal{V}^{(p)}$ at step $t$ is the smallest set satisfying:
\begin{equation}
  \sum_{v \in \mathcal{V}^{(p)}} P(v \mid x, s_{<t}) \ge p
\end{equation}
Tokens outside $\mathcal{V}^{(p)}$ are assigned probability 0 before final selection.

\subsection{Architecture Flow Diagram}

\begin{figure}[h]
\centering
\begin{tikzpicture}[node distance=0.6cm and 0.9cm]
  \node[datblk] (q) {Question $x$};
  \node[grayblk, right=of q] (cot) {Reasoning\\Chain $c_1..c_M$};
  \node[trnblk, right=of cot] (ans) {Final\\Answer $y$};
  \node[lossblk, below=0.8cm of cot] (lc) {Chain Loss\\$\mathcal{L}_c$};
  \node[lossblk, below=0.8cm of ans] (la) {Answer Loss\\$\mathcal{L}_y$};
  \node[mrgblk, below=0.8cm of $(lc)!0.5!(la)$] (total) {Total Loss\\$\mathcal{L}_{\text{CoT}}$};

  \draw[arr] (q)--(cot);
  \draw[arr] (cot)--(ans);
  \draw[arr] (cot)--(lc);
  \draw[arr] (ans)--(la);
  \draw[arr] (lc)--(total);
  \draw[arr] (la)--(total);

  \draw[losscolor,thick,decorate,decoration={brace,amplitude=6pt,raise=4pt}]
    (cot.north west) -- (ans.north east) node[midway,above=10pt,font=\tiny]{loss applied to entire output};
\end{tikzpicture}
\caption{CoT SFT: loss is computed over both the reasoning chain and the final answer.}
\end{figure}

\subsection{State Transition Table}
\begin{table}[h]\centering
\begin{tabular}{lll}\toprule
\textbf{Property} & \textbf{Before} & \textbf{After}\\\midrule
Output format & Direct answer & Reasoning chain + answer\\
Training signal per example & $|y|$ tokens & $|c| + |y|$ tokens\\
Multi-step reasoning & Weak & Strong\\
Data annotation cost & Low & High (chain annotation)\\
Accuracy on complex tasks & Low & Substantially improved\\\bottomrule
\end{tabular}
\caption{State properties before and after Chain-of-Thought SFT.}
\end{table}

\newpage
%% ------------------------------------------------------------
\section{Step-by-Step SFT}

\subsection{Mathematical Foundation}
Step-by-Step SFT formats the reasoning chain into discrete steps separated by a boundary token $T_{\text{step}} = \text{\texttt{[STEP]}}$. Let the response contain $S$ steps: $s = (s_1, T_{\text{step}}, s_2, T_{\text{step}}, \dots, s_S, T_{\text{step}})$. The step-level loss is:
\begin{equation}
  \mathcal{L}_{\text{SBS}}(\theta) = -\sum_{k=1}^S \sum_{t=1}^{|s_k|} \log P_\theta(s_{k,t} \mid x, s_{<k}, s_{k,<t}) - \sum_{k=1}^S \log P_\theta(T_{\text{step}} \mid x, s_{\le k})
\end{equation}
We define the average step perplexity $\text{PPL}_k$ as:
\begin{equation}
  \text{PPL}_k = \exp\left( -\frac{1}{|s_k|} \sum_{t=1}^{|s_k|} \log P_\theta(s_{k,t} \mid x, s_{<k}, s_{k,<t}) \right)
\end{equation}
Steps with high perplexity ($\text{PPL}_k > \tau$, where $\tau$ is a threshold) indicate low model confidence and can be flagged for step-level revision during generation.
The model is trained to minimize the combined cross-entropy loss over all steps.

\subsection{Architecture Flow Diagram}

\begin{figure}[h]
\centering
\begin{tikzpicture}[node distance=0.5cm and 0.8cm]
  \node[datblk,minimum width=1.8cm] (inp) {Input $x$};
  \node[trnblk,right=0.7cm of inp,minimum width=1.8cm] (s1) {Step 1\\$s_1$};
  \node[trnblk,right=0.6cm of s1,minimum width=1.8cm] (s2) {Step 2\\$s_2$};
  \node[grayblk,right=0.6cm of s2,minimum width=0.6cm] (dots) {$\cdots$};
  \node[trnblk,right=0.6cm of dots,minimum width=1.8cm] (sn) {Step $N$\\$s_N$};
  \node[mrgblk,right=0.7cm of sn,minimum width=1.8cm] (out) {Answer\\$y$};

  \foreach \a/\b in {inp/s1, s1/s2, s2/dots, dots/sn, sn/out}
    \draw[arr] (\a)--(\b);

  \node[losscolor,font=\tiny] at ($(s1.south)+(0,-0.45)$) {$\mathcal{L}_1$};
  \node[losscolor,font=\tiny] at ($(s2.south)+(0,-0.45)$) {$\mathcal{L}_2$};
  \node[losscolor,font=\tiny] at ($(sn.south)+(0,-0.45)$) {$\mathcal{L}_N$};
  \node[losscolor,font=\tiny] at ($(out.south)+(0,-0.45)$) {$\mathcal{L}_y$};

  \foreach \n in {s1,s2,sn,out}
    \draw[garr] (\n.south) -- ++(0,-0.35);
\end{tikzpicture}
\caption{Step-by-Step SFT: each numbered step and the final answer contribute to the total loss.}
\end{figure}

\subsection{State Transition Table}
\begin{table}[h]\centering
\begin{tabular}{lll}\toprule
\textbf{Property} & \textbf{Before} & \textbf{After}\\\midrule
Solution structure & Unstructured & Numbered procedural steps\\
Step delimiter tokens & Not in vocabulary & Learned as first-class tokens\\
Interpretability & Low & Step-level traceability\\
Annotation format & Raw answer & $(x, [s_1,...,s_N], y)$\\
Task decomposition ability & Implicit & Explicit\\\bottomrule
\end{tabular}
\caption{State properties before and after Step-by-Step SFT.}
\end{table}

\newpage
%% ------------------------------------------------------------
\section{Process Supervision}

\subsection{Mathematical Foundation}
Process Supervision trains a Process Reward Model (PRM) using step-level binary correctness labels. Let $x$ be the problem description, and $s_1, \dots, s_S$ be the generated steps. Human or model-based annotations assign correctness labels $r_k \in \{0, 1\}$ for each step $s_k$.
The PRM model outputs correctness probabilities by applying a logistic sigmoid function to the step boundary token representation:
\begin{equation}
  \hat{r}_k = \sigma\left( \text{PRM}_\phi(h_{t_k}) \right) = \frac{1}{1 + \exp\left( -\text{PRM}_\phi(h_{t_k}) \right)}
\end{equation}
where $h_{t_k}$ is the transformer hidden output at index $t_k$ representing the step boundary token. The PRM parameters $\phi$ are optimized using binary cross entropy:
\begin{equation}
  \mathcal{L}_{\text{PRM}}(\phi) = -\sum_{k=1}^S \left[ r_k \log \hat{r}_k + (1-r_k) \log(1 - \hat{r}_k) \right]
\end{equation}
During inference, candidate sequences are generated. The total sequence correctness score is calculated as the product of step correctness probabilities:
\begin{equation}
  \text{Score}(s_{\le S}) = \prod_{k=1}^S \hat{r}_k
\end{equation}
This score is used to rank candidate rollouts in Best-of-$N$ sampling or guide selection during Monte Carlo Tree Search (MCTS).

\subsection{Architecture Flow Diagram}

\begin{figure}[h]
\centering
\begin{tikzpicture}[node distance=0.6cm and 1.0cm]
  \node[datblk] (sol) {Solution\\Steps $s_{1..N}$};
  \node[grayblk, right=of sol] (prm) {Process Reward\\Annotator};
  \node[trnblk, right=0.6cm of prm, minimum height=2.2cm, minimum width=2.2cm] (labels) {Step\\Labels\\$r_1,..,r_N$};
  \node[lossblk, right=of labels] (loss) {Weighted\\SFT Loss};
  \node[mrgblk, below=0.9cm of loss] (out) {Process-Supervised\\Model};

  \node[green!60!black,draw,rounded corners,font=\tiny,minimum width=1.5cm] at ($(labels.east)+(0,0.5)$) {$r_n{=}{+}1$: train};
  \node[red!70!black,draw,rounded corners,font=\tiny,minimum width=1.5cm] at ($(labels.east)+(0,-0.5)$) {$r_n{=}{-}1$: skip};

  \draw[arr] (sol)--(prm);
  \draw[arr] (prm)--(labels);
  \draw[arr] (labels)--(loss);
  \draw[arr] (loss)--(out);
\end{tikzpicture}
\caption{Process supervision: a step annotator labels each reasoning step; only correct steps contribute to the SFT loss.}
\end{figure}

\subsection{State Transition Table}
\begin{table}[h]\centering
\begin{tabular}{lll}\toprule
\textbf{Property} & \textbf{Before} & \textbf{After}\\\midrule
Supervision granularity & Final answer only & Per-step labels\\
Error signal & Outcome-level & Step-level\\
Error localisation & None & Identifies which step failed\\
Annotation cost & Low & High (step-level labelling)\\
Reasoning quality & Outcome-driven & Step-correctness driven\\\bottomrule
\end{tabular}
\caption{State properties before and after process supervision.}
\end{table}

\newpage
%% ============================================================
\part{Parameter-Efficient Fine-Tuning (PEFT)}
%% ============================================================

\section{Low-Rank Adaptation (LoRA)}

\subsection{Mathematical Foundation}
LoRA represents the weight update matrix $\Delta W \in \mathbb{R}^{k \times d}$ for a linear projection layer as the product of two low-rank matrices $B \in \mathbb{R}^{k \times r}$ and $A \in \mathbb{R}^{r \times d}$, where the rank satisfies $r \ll \min(d, k)$:
\begin{equation}
  h = W_0 x + \Delta W x = W_0 x + \frac{\alpha}{r} B A x
\end{equation}
where $\alpha$ is a constant scaling parameter. The base weight matrix $W_0 \in \mathbb{R}^{k \times d}$ is frozen; only $B$ and $A$ are updated. Under a loss $\mathcal{L}$, the gradient propagation rules using the chain rule are derived as follows:
Let $z = A x \in \mathbb{R}^r$ be the intermediate bottleneck activation. Then $h = W_0 x + \frac{\alpha}{r} B z$.
The gradient with respect to output activation is $\nabla_h \mathcal{L} \in \mathbb{R}^{k \times 1}$ (represented as a column vector).
The gradient with respect to the intermediate bottleneck representation $z$:
\begin{equation}
  \nabla_z \mathcal{L} = \frac{\alpha}{r} B^\top \nabla_h \mathcal{L} \in \mathbb{R}^{r \times 1}
\end{equation}
The gradient with respect to the matrix $B$:
\begin{equation}
  \nabla_B \mathcal{L} = \frac{\alpha}{r} \nabla_h \mathcal{L} z^\top = \frac{\alpha}{r} \nabla_h \mathcal{L} x^\top A^\top \in \mathbb{R}^{k \times r}
\end{equation}
The gradient with respect to the matrix $A$:
\begin{equation}
  \nabla_A \mathcal{L} = \nabla_z \mathcal{L} x^\top = \frac{\alpha}{r} B^\top \nabla_h \mathcal{L} x^\top \in \mathbb{R}^{r \times d}
\end{equation}
To eliminate latency at deployment, the weights are merged directly into the base model:
\begin{equation}
  W_{\text{merged}} = W_0 + \frac{\alpha}{r} B A
\end{equation}

\subsection{Architecture Flow Diagram}

\begin{figure}[h]
\centering
\begin{tikzpicture}[node distance=0.7cm and 1.1cm]
  \node[neuron] (x1) at (0,1.5) {$x_1$};
  \node[neuron] (x2) at (0,0.5) {$x_2$};
  \node[neuron] (xd) at (0,-0.8) {$x_d$};
  \node[font=\tiny] at (0,-0.15) {$\vdots$};

  \node[frzblk, minimum width=2.8cm, minimum height=2.0cm] (W0) at (3.8,1.5) {Frozen\\$W_0 \in \mathbb{R}^{k\times d}$};
  \node[bneck] (z1) at (3.8,-0.2) {$z_1$};
  \node[bneck] (zr) at (3.8,-1.2) {$z_r$};
  \node[font=\tiny] at (3.8,-0.7) {$\vdots$};

  \node[outnrn] (h1) at (7.6,1.5) {$h_1$};
  \node[outnrn] (h2) at (7.6,0.5) {$h_2$};
  \node[outnrn] (hk) at (7.6,-0.8) {$h_k$};
  \node[font=\tiny] at (7.6,-0.15) {$\vdots$};

  \draw[arr] (x1.east) -- ($(W0.west)+(0,0.5)$);
  \draw[arr] (x2.east) -- ($(W0.west)+(0,0)$);
  \draw[arr] (xd.east) -- ($(W0.west)+(0,-0.5)$);
  \foreach \s in {x1,x2,xd} \foreach \t in {z1,zr}
    \draw[arr,green!70!black,thin] (\s.east)--(\t.west);
  \foreach \s in {z1,zr} \foreach \t in {h1,h2,hk}
    \draw[arr,green!70!black,thin] (\s.east)--(\t.west);
  \draw[arr] ($(W0.east)+(0,0.5)$)--(h1.west);
  \draw[arr] ($(W0.east)+(0,0.0)$)--(h2.west);
  \draw[arr] ($(W0.east)+(0,-0.5)$)--(hk.west);

  \node[green!60!black,font=\tiny] at (1.9,-0.9) {$A$ (train)};
  \node[green!60!black,font=\tiny] at (5.7,-0.9) {$B$ (train)};
  \draw[garr] (h1.east)--++(0.6,0) node[right,red,font=\tiny]{$g_1$};
  \draw[garr] (hk.east)--++(0.6,0) node[right,red,font=\tiny]{$g_k$};
\end{tikzpicture}
\caption{LoRA: input flows through frozen $W_0$ and trainable low-rank path $A \to B$; outputs are summed.}
\end{figure}

\subsection{State Transition Table}
\begin{table}[h]\centering
\begin{tabular}{lll}\toprule
\textbf{Property} & \textbf{Before} & \textbf{After}\\\midrule
Base weights $W_0$ & requires\_grad = True & Frozen (False)\\
Trainable params & $k \times d$ & $r(k+d) \ll kd$\\
Matrix $B$ & Does not exist & $\mathbb{R}^{k\times r}$, init $= 0$\\
Matrix $A$ & Does not exist & $\mathbb{R}^{r\times d}$, init $\sim \mathcal{N}$\\
Inference overhead & None & Extra BA multiply (or merged)\\\bottomrule
\end{tabular}
\caption{State properties before and after LoRA injection.}
\end{table}

\newpage
%% ------------------------------------------------------------
\section{Quantized LoRA (QLoRA)}

\subsection{Mathematical Foundation}
QLoRA quantizes the base weight $W_0$ to 4-bit NormalFloat (NF4) and uses Double Quantization (DQ) to minimize memory requirements.
NF4 divides the standard normal distribution $\mathcal{N}(0, 1)$ into $2^k=16$ equal-area intervals. The quantization levels $q_i$ are calculated using the standard normal quantile function $Q_X$:
\begin{equation}
  q_i = \frac{1}{2} \left( Q_X\left(\frac{i}{2^k}\right) + Q_X\left(\frac{i+1}{2^k}\right) \right)
\end{equation}
The base weights are normalized to $[-1, 1]$ using block-wise scaling. Let the block size be $B=64$. For a block $W_b$, the scale is $c_b = \max(|W_b|)$. The quantized values are:
\begin{equation}
  q_t = \arg\min_{i \in \{0, \dots, 15\}} \left| \frac{W_{b,t}}{c_b} - q_i \right|
\end{equation}
Double Quantization (DQ) quantizes the scaling factors $c_b$ themselves to 8-bit floats ($c_b^{\text{FP8}}$) with a second scale block size of 256, reducing scale memory overhead from $32/64 = 0.5$ bits/weight to $8/64 + 32/(64 \times 256) \approx 0.127$ bits/weight.
During backpropagation, the NF4 weights are dequantized to FP16/BF16 to compute outputs, and gradients are computed only for the low-rank parameters:
\begin{equation}
  h = \text{dequant}(c_b, W^{\text{NF4}}) x + \frac{\alpha}{r} B A x
\end{equation}

\subsection{Architecture Flow Diagram}

\begin{figure}[h]
\centering
\begin{tikzpicture}[node distance=0.7cm and 1.2cm]
  \node[datblk] (x) {Input\\$x$ (BF16)};
  \node[frzblk, right=of x] (nf4) {$W_0^{\text{NF4}}$\\4-bit frozen};
  \node[grayblk, right=of nf4] (deq) {Dequantise\\to BF16};
  \node[sumnode, right=of deq] (sum) {$\oplus$};
  \node[trnblk, above=0.8cm of sum] (lora) {LoRA\\$BA$ (BF16)};
  \node[mrgblk, right=of sum] (out) {Output\\$h$ (BF16)};

  \draw[arr] (x)--(nf4);
  \draw[arr] (nf4)--(deq);
  \draw[arr] (deq)--(sum);
  \draw[arr] (lora)--(sum);
  \draw[arr] (x.north) |- (lora.west);
  \draw[arr] (sum)--(out);

  \node[red,font=\tiny,text width=2.5cm,align=center,below=0.3cm of nf4] {4-bit: $\sim$4 GB\\vs BF16: $\sim$16 GB};
\end{tikzpicture}
\caption{QLoRA: 4-bit frozen base dequantised at runtime; LoRA adapters trained in BF16.}
\end{figure}

\subsection{State Transition Table}
\begin{table}[h]\centering
\begin{tabular}{lll}\toprule
\textbf{Property} & \textbf{Before} & \textbf{After}\\\midrule
$W_0$ precision & BF16 / FP16 & NF4 (4-bit)\\
GPU memory for 70B model & ${\sim}140$ GB & ${\sim}35$ GB\\
LoRA adapter precision & N/A & BF16\\
Training fidelity & Full & Near-full (straight-through)\\
Inference latency & Baseline & $+$dequant overhead\\\bottomrule
\end{tabular}
\caption{State properties before and after QLoRA.}
\end{table}

\newpage
%% ------------------------------------------------------------
\section{Adaptive LoRA (AdaLoRA)}

\subsection{Mathematical Foundation}
AdaLoRA parameterizes incremental weight updates using a singular value decomposition (SVD) representation to dynamically adjust the adapter rank during training:
\begin{equation}
  \Delta W = P \Lambda Q
\end{equation}
where $P \in \mathbb{R}^{d \times r}$ and $Q \in \mathbb{R}^{r \times k}$ are left and right singular vectors, and $\Lambda = \text{diag}(\lambda_1, \dots, \lambda_r) \in \mathbb{R}^{r \times r}$ contains singular values.
To prevent SVD divergence and maintain parameter stability, we enforce orthogonality constraints:
\begin{equation}
  \mathcal{L}_{\text{orth}}(P, Q) = \gamma \left( \|P^\top P - I\|_F^2 + \|Q Q^\top - I\|_F^2 \right)
\end{equation}
We estimate the training importance of each singular triplet $(P_{*,i}, \lambda_i, Q_{i,*})$ using gradient-weight product magnitude:
\begin{equation}
  S_i = s(\lambda_i) + \frac{1}{d} \sum_{j=1}^d s(P_{j,i}) + \frac{1}{k} \sum_{j=1}^k s(Q_{i,j})
\end{equation}
where $s(\theta) = |\theta \cdot \nabla_\theta \mathcal{L}|$. At designated intervals, triplets with the lowest importance scores $S_i$ are pruned by zeroing out their singular values $\lambda_i$, dynamically reducing the adapter rank budget.

\subsection{Architecture Flow Diagram}

\begin{figure}[h]
\centering
\begin{tikzpicture}[node distance=0.7cm and 1.1cm]
  \node[datblk] (x) {Input $x$};
  \node[frzblk, right=1.5cm of x] (W0) {Frozen $W_0$};
  
  % Parallel SVD Adapter path
  \node[trnblk, above=1.0cm of W0] (Q) {$Q$ (down-proj)};
  \node[trnblk, right=1.2cm of Q] (L) {$\Lambda$ (diag)};
  \node[trnblk, right=1.2cm of L] (P) {$P$ (up-proj)};
  
  \node[grayblk, above=0.6cm of L] (prune) {Rank Pruner};
  \node[sumnode, right=1.5cm of W0] (sum) {$\oplus$};
  \node[mrgblk, right=1.5cm of sum] (out) {Output $h$};

  % Connections
  \draw[arr] (x)--(W0);
  \draw[arr] (x.north) |- (Q.west);
  \draw[arr] (Q)--(L);
  \draw[arr] (L)--(P);
  \draw[arr] (P.east) -| (sum.north);
  \draw[arr] (W0)--(sum);
  \draw[arr] (sum)--(out);
  
  % Pruning monitor line
  \draw[arr,dashed,red] (prune)--node[right,font=\tiny]{prunes $\lambda_i$}(L);
  \node[red,font=\tiny,align=center] at ($(prune.north)+(0,0.25)$) {$S_i{<}\gamma \Rightarrow \lambda_i{=}0$};
\end{tikzpicture}
\caption{AdaLoRA: SVD-decomposed update with importance-based rank pruning per layer.}
\end{figure}

\subsection{State Transition Table}
\begin{table}[h]\centering
\begin{tabular}{lll}\toprule
\textbf{Property} & \textbf{Before} & \textbf{After}\\\midrule
Rank per layer & Fixed $r$ & Adaptive (varies by importance)\\
Update structure & $BA$ product & $P\Lambda Q$ SVD triplet\\
Parameter budget & Fixed & Redistributed across layers\\
Low-importance ranks & Included & Pruned ($\lambda_i = 0$)\\
Training stability & Standard & Constrained by SVD orthogonality\\\bottomrule
\end{tabular}
\caption{State properties before and after AdaLoRA.}
\end{table}

\newpage
%% ------------------------------------------------------------
\section{Weight-Decomposed LoRA (DoRA)}

\subsection{Mathematical Foundation}
DoRA decomposes the weight matrix $W \in \mathbb{R}^{k \times d}$ into its magnitude vector $m \in \mathbb{R}^k$ and direction matrix $V \in \mathbb{R}^{k \times d}$:
\begin{equation}
  W = m \cdot \frac{V}{\|V\|_r} = m \cdot \frac{W_0 + \Delta W}{\|W_0 + \Delta W\|_r}
\end{equation}
where $\|\cdot\|_r$ is the L2 norm computed along each row (corresponding to output channels). In parameter-efficient training, we freeze the direction matrix to the base weight $V = W_0$ and introduce low-rank directional updates using adapters $\Delta V = \frac{\alpha}{r} B A$ (where $A \in \mathbb{R}^{r \times d}$ and $B \in \mathbb{R}^{k \times r}$):
\begin{equation}
  W = m \cdot \frac{W_0 + \frac{\alpha}{r} B A}{\|W_0 + \frac{\alpha}{r} B A\|_r}
\end{equation}
Only the magnitude vector $m$ and low-rank directional matrices $B$ and $A$ are trainable.
The gradient flow for the directional components is:
\begin{equation}
  \nabla_V \mathcal{L} = \frac{m}{\|V\|_c} \left( \nabla_W \mathcal{L} - \frac{V}{\|V\|_c^2} \odot \left( V^\top \nabla_W \mathcal{L} \right) \right)
\end{equation}
This mathematically separates the updates to magnitude and direction, mimicking full parameter fine-tuning behavior while retaining low parameter overhead.

\subsection{Architecture Flow Diagram}

\begin{figure}[h]
\centering
\begin{tikzpicture}[node distance=0.7cm and 1.1cm]
  \node[frzblk] (W0) {$W_0$ (frozen)};
  \node[trnblk, above=1.0cm of W0] (ba) {$BA$ (LoRA)};
  
  \node[sumnode, right=1.2cm of W0] (sumV) {$\oplus$};
  
  \node[grayblk, right=1.0cm of sumV] (norm) {Norm $\|\cdot\|_r$};
  \node[sumnode, right=1.2cm of norm] (div) {$\div$};
  
  \node[trnblk, above=1.0cm of div] (m) {Magnitude $m$};
  \node[sumnode, right=1.2cm of div] (mul) {$\otimes$};
  \node[mrgblk, right=1.2cm of mul] (out) {Weight $W'$};

  % Connections
  \draw[arr] (W0)--(sumV);
  \draw[arr] (ba.east) -| (sumV.north);
  \draw[arr] (sumV)--(norm);
  \draw[arr] (sumV.east) -- ++(0.2,0) |- ($(div.west)+(0,0.2)$);
  \draw[arr] (norm)--(div);
  \draw[arr] (div)--(mul);
  \draw[arr] (m)--(mul);
  \draw[arr] (mul)--(out);

  \node[font=\tiny,below=0.3cm of norm] {direction component};
  \node[font=\tiny,above=0.3cm of m] {amplitude scaling};
\end{tikzpicture}
\caption{DoRA: weight decomposed into direction (LoRA-adapted) and magnitude (trainable scalar).}
\end{figure}

\subsection{State Transition Table}
\begin{table}[h]\centering
\begin{tabular}{lll}\toprule
\textbf{Property} & \textbf{Before (LoRA)} & \textbf{After (DoRA)}\\\midrule
Magnitude component & Implicit in $W_0$ & Explicit trainable $m$\\
Direction component & $W_0 + BA$ unnorm & Column-normalised\\
Full FT similarity & Lower & Higher\\
Trainable params & $r(d+k)$ & $r(d+k) + k$\\
Catastrophic forgetting & Low & Lower\\\bottomrule
\end{tabular}
\caption{State properties comparing LoRA and DoRA.}
\end{table}

\newpage
%% ------------------------------------------------------------
\section{IA$^3$ (Infused Adapter by Inhibiting and Amplifying Inner Activations)}

\subsection{Mathematical Foundation}
IA$^3$ introduces element-wise scaling vectors to directly modify inner activations in attention keys, values, and intermediate feed-forward layers. Let the scaling vectors be $l_K \in \mathbb{R}^d$, $l_V \in \mathbb{R}^d$, and $l_{FFN} \in \mathbb{R}^{d_{FFN}}$.
The attention layer activations are computed as:
\begin{align}
  \tilde{K} &= K \odot l_K = K \cdot \text{diag}(l_K) \\
  \tilde{V} &= V \odot l_V = V \cdot \text{diag}(l_V) \\
  \text{Attention}(Q, \tilde{K}, \tilde{V}) &= \text{softmax}\left( \frac{Q \tilde{K}^\top}{\sqrt{d_k}} \right) \tilde{V}
\end{align}
In the feed-forward network (FFN), let $W_1$ and $W_2$ be the weight matrices. The intermediate activations are scaled before the output projection:
\begin{equation}
  \text{FFN}(x) = \sigma\bigl(x W_1 \odot l_{FFN}\bigr) W_2 = \sigma\bigl(x W_1 \cdot \text{diag}(l_{FFN})\bigr) W_2
\end{equation}
The base model weights $W_Q, W_K, W_V, W_1, W_2$ are frozen. This diagonal matrix transformation scales the columns of the projection matrices with minimal parameter updates.

\subsection{Architecture Flow Diagram}

\begin{figure}[h]
\centering
\begin{tikzpicture}[node distance=0.6cm and 1.0cm]
  \node[datblk] (x) {Input $x$};
  
  % Attention Block
  \node[frzblk, right=1.0cm of x] (qkv) {Compute\\$Q, K, V$};
  \node[trnblk, above right=0.3cm and 0.8cm of qkv] (lk) {$l_K \odot K$};
  \node[trnblk, below right=0.3cm and 0.8cm of qkv] (lv) {$l_V \odot V$};
  \node[frzblk, right=2.8cm of qkv] (attn) {Attention\\Softmax};
  
  % FFN Block
  \node[frzblk, right=1.0cm of attn] (w1) {FFN $W_1$};
  \node[trnblk, right=0.8cm of w1] (lffn) {$l_{FFN} \odot$};
  \node[frzblk, right=0.8cm of lffn] (w2) {FFN $W_2$};
  \node[mrgblk, right=1.0cm of w2] (out) {Output};

  % Connections
  \draw[arr] (x)--(qkv);
  \draw[arr] (qkv.east) -- ++(0.2,0) |- (lk.west);
  \draw[arr] (qkv.east) -- ++(0.2,0) |- (lv.west);
  \draw[arr] (qkv.east) -- ++(0.2,0) |- ($(attn.west)+(0,0.4)$); % Q direct
  \draw[arr] (lk.east) -| ($(attn.west)+(0,0.15)$);
  \draw[arr] (lv.east) -| ($(attn.west)+(0,-0.15)$);
  
  \draw[arr] (attn)--(w1);
  \draw[arr] (w1)--(lffn);
  \draw[arr] (lffn)--(w2);
  \draw[arr] (w2)--(out);
\end{tikzpicture}
\caption{IA$^3$: three tiny scaling vectors rescale keys, values, and FFN intermediate activations.}
\end{figure}

\subsection{State Transition Table}
\begin{table}[h]\centering
\begin{tabular}{lll}\toprule
\textbf{Property} & \textbf{Before} & \textbf{After}\\\midrule
Trainable params & 0 (frozen) & $3dL$ scalars\\
Init value of $l$ vectors & N/A & $\mathbf{1}$ (identity behaviour)\\
K, V, FFN activations & Unchanged & Element-wise scaled\\
Param overhead vs LoRA & N/A & ${\sim}10\times$ fewer params\\
Task adaptation strength & None & Moderate\\\bottomrule
\end{tabular}
\caption{State properties before and after IA$^3$ injection.}
\end{table}

\newpage
%% ------------------------------------------------------------
\section{Adapter Tuning}

\subsection{Mathematical Foundation}
Adapter tuning inserts small bottleneck layers into the transformer block, typically after attention and feed-forward modules. Let $h \in \mathbb{R}^d$ be the input hidden state.
In the Pfeiffer architecture, a single adapter is inserted per block. In the Houlsby architecture, two adapters are inserted per block. The mathematical transformation for an adapter block is:
\begin{equation}
  h' = h + \sigma(h W_{down}) W_{up}
\end{equation}
where $W_{down} \in \mathbb{R}^{d \times r}$ projects to a bottleneck dimension $r \ll d$, $\sigma$ is an activation function (typically GELU), and $W_{up} \in \mathbb{R}^{r \times d}$ projects back to the original dimension.
To ensure the adapter behaves as an identity mapping at the start of training, $W_{up}$ is initialized to zero:
\begin{equation}
  h' \bigr|_{t=0} = h + \sigma(h W_{down}) \cdot 0 = h
\end{equation}

\subsection{Architecture Flow Diagram}

\begin{figure}[h]
\centering
\begin{tikzpicture}[node distance=0.5cm and 0.9cm]
  \node[datblk,minimum width=2.2cm] (h) {Hidden $h$\\(from layer)};
  \node[trnblk,right=of h,minimum width=2.0cm] (dn) {Down-proj\\$W_{\text{dn}}$ ($d{\to}r$)};
  \node[grayblk,right=of dn,minimum width=1.8cm] (act) {ReLU};
  \node[trnblk,right=of act,minimum width=2.0cm] (up) {Up-proj\\$W_{\text{up}}$ ($r{\to}d$)};
  \node[sumnode,right=of up] (add) {$+$};
  \node[mrgblk,right=of add,minimum width=2.0cm] (out) {Adapter\\Output $h'$};

  \draw[arr] (h)--(dn);
  \draw[arr] (dn)--(act);
  \draw[arr] (act)--(up);
  \draw[arr] (up)--(add);
  \draw[arr] (h.south) -- ++(0,-0.6) -| (add.south) node[midway,below,font=\tiny]{skip connection};
  \draw[arr] (add)--(out);

  \node[red,font=\tiny,below=0.2cm of dn] {bottleneck $r \ll d$};
\end{tikzpicture}
\caption{Adapter module: down-project to rank $r$, apply nonlinearity, up-project back, plus skip connection.}
\end{figure}

\subsection{State Transition Table}
\begin{table}[h]\centering
\begin{tabular}{lll}\toprule
\textbf{Property} & \textbf{Before} & \textbf{After}\\\midrule
Sub-layer output path & Direct pass-through & $+$adapter bottleneck\\
Transformer weights & Trainable & Frozen\\
New params per layer & 0 & $2dr + 2d$\\
Adapter init & N/A & Near-identity ($W_{\text{dn}} \sim \mathcal{N}$)\\
Inference latency & Baseline & $+$two extra matmuls per layer\\\bottomrule
\end{tabular}
\caption{State properties before and after adapter insertion.}
\end{table}

\newpage
%% ------------------------------------------------------------
\section{Prefix Tuning}

\subsection{Mathematical Foundation}
Prefix Tuning prepends continuous task-specific virtual prefix vectors $P_K, P_V \in \mathbb{R}^{l \times d}$ directly to the Key and Value matrices at every transformer layer. Let the original Key and Value matrices for sequence length $T$ be $K, V \in \mathbb{R}^{T \times d}$. The modified Key and Value matrices become:
\begin{align}
  \tilde{K} &= [P_K ; K] \in \mathbb{R}^{(l+T) \times d} \\
  \tilde{V} &= [P_V ; V] \in \mathbb{R}^{(l+T) \times d}
\end{align}
To stabilize optimization, during training the prefix is reparameterized using a Multilayer Perceptron (MLP) over a small source embedding $E_k$:
\begin{equation}
  P_K = \text{MLP}_K(E_k), \quad P_V = \text{MLP}_V(E_k)
\end{equation}
where $E_k \in \mathbb{R}^{l \times d_{\text{mid}}}$. The MLP is discarded at inference, and only the generated static prefix matrices $P_K, P_V$ are saved and used.
The attention output is computed as:
\begin{equation}
  \text{Attention}(Q, \tilde{K}, \tilde{V}) = \text{softmax}\left( \frac{Q [P_K ; K]^\top}{\sqrt{d_k}} \right) [P_V ; V]
\end{equation}

\subsection{Architecture Flow Diagram}

\begin{figure}[h]
\centering
\begin{tikzpicture}[node distance=0.6cm and 1.0cm]
  \node[trnblk,minimum width=2.5cm] (prefix) {Trainable Prefix\\$P_k, P_v$ (per layer)};
  \node[datblk, below=1.2cm of prefix] (tokens) {Input Token\\Keys \& Values};
  \node[grayblk, right=1.5cm of $(prefix)!0.5!(tokens)$] (concat) {Concat\\$[P_k;K],\,[P_v;V]$};
  \node[frzblk, right=of concat] (attn) {Frozen\\Attention};
  \node[mrgblk, right=of attn] (out) {Contextual\\Output};

  \draw[arr] (prefix.east) -| (concat.north);
  \draw[arr] (tokens.east) -| (concat.south);
  \draw[arr] (concat)--(attn);
  \draw[arr] (attn)--(out);

  \node[font=\tiny,green!60!black] at ($(prefix.south)+(0,-0.2)$) {only these trained};
  \node[font=\tiny,blue!60!black] at ($(attn.south)+(0,-0.2)$) {all frozen};
\end{tikzpicture}
\caption{Prefix Tuning: trainable prefix vectors extend the key/value context at every attention layer.}
\end{figure}

\subsection{State Transition Table}
\begin{table}[h]\centering
\begin{tabular}{lll}\toprule
\textbf{Property} & \textbf{Before} & \textbf{After}\\\midrule
Attention context length & $T$ tokens & $T + l$ tokens (prefix prepended)\\
Transformer params & Trainable & Fully frozen\\
Trainable params & 0 & $2 \times l \times d_h \times L$\\
Prefix init & N/A & Via MLP reparameterisation\\
Task switch at inference & Reload model & Swap prefix only\\\bottomrule
\end{tabular}
\caption{State properties before and after prefix tuning.}
\end{table}

\newpage
%% ------------------------------------------------------------
\section{Prompt Tuning}

\subsection{Mathematical Foundation}
Prompt tuning prepends $p$ learnable embedding vectors $P \in \mathbb{R}^{p \times d_{emb}}$ to the input word representations:
\begin{equation}
  \tilde{X} = [P ; X_e] \in \mathbb{R}^{(p + T) \times d_{emb}}
\end{equation}
where $X_e \in \mathbb{R}^{T \times d_{emb}}$ is the original sequence word embedding matrix. The entire transformer backbone parameters $\theta_{base}$ remain frozen, and gradients are backpropagated only to the prompt parameters $P$:
\begin{equation}
  \nabla_P \mathcal{L} = \frac{\partial \mathcal{L}}{\partial \tilde{X}_{1..p}}
\end{equation}
During backpropagation, the gradients are passed through all layers of the frozen transformer and update only the prompt embeddings.

\subsection{Architecture Flow Diagram}

\begin{figure}[h]
\centering
\begin{tikzpicture}[node distance=0.7cm and 1.0cm]
  \node[trnblk,minimum width=2.2cm] (soft) {Soft Prompt\\$P \in \mathbb{R}^{k\times d}$\\(trainable)};
  \node[datblk, right=of soft] (emb) {Input\\Embeddings\\$E(x_{1..T})$};
  \node[grayblk, right=of emb] (concat) {Concat\\$[\tilde{x}]$};
  \node[frzblk, right=of concat] (lm) {Frozen\\Transformer\\$f_\theta$};
  \node[lossblk, right=of lm] (loss) {SFT Loss\\on $y$};

  \draw[arr] (soft.east) -| (concat.north);
  \draw[arr] (emb)--(concat);
  \draw[arr] (concat)--(lm);
  \draw[arr] (lm)--(loss);

  \node[font=\tiny,green!60!black,align=center,below=0.25cm of soft] {only $P$\\is trained};
\end{tikzpicture}
\caption{Prompt Tuning: soft tokens prepended at the embedding layer; the full transformer is frozen.}
\end{figure}

\subsection{State Transition Table}
\begin{table}[h]\centering
\begin{tabular}{lll}\toprule
\textbf{Property} & \textbf{Before} & \textbf{After}\\\midrule
Trainable params & Full model & $k \times d_{\text{emb}}$ only\\
Prompt format & Discrete hand-crafted & Continuous soft vectors\\
Transformer weights & All trainable & All frozen\\
Task switching & Reload fine-tuned model & Swap $P$ only\\
Sensitivity to init & N/A & Strong (init from class labels helps)\\\bottomrule
\end{tabular}
\caption{State properties before and after prompt tuning.}
\end{table}

\newpage
%% ------------------------------------------------------------
\section{P-Tuning v2}

\subsection{Mathematical Foundation}
P-Tuning v2 extends prompt tuning by prepending learnable virtual tokens $P^{(l)} \in \mathbb{R}^{p \times 2 d_{attn}}$ directly to the attention key and value states at every layer $l \in [1, L]$ of the network. Unlike prefix tuning, P-Tuning v2 removes the MLP reparameterization network, directly optimizing the prefix parameters $P^{(l)}$:
\begin{align}
  \tilde{K}^{(l)} &= [P^{(l)}_K ; K^{(l)}] \\
  \tilde{V}^{(l)} &= [P^{(l)}_V ; V^{(l)}]
\end{align}
This deep prefix injection resolves the capacity limitation of shallow prompt tuning, allowing small models ($<10\text{B}$ parameters) to achieve performance comparable to full fine-tuning.

\subsection{Architecture Flow Diagram}

\begin{figure}[h]
\centering
\begin{tikzpicture}[node distance=0.5cm and 0.9cm]
  \node[datblk,minimum width=2.0cm] (emb) {Embedding\\Layer};
  \node[frzblk, right=0.9cm of emb,minimum width=2.0cm] (l1) {Layer 1\\(Attn + FFN)};
  \node[frzblk, right=0.9cm of l1,minimum width=2.0cm] (l2) {Layer 2\\(Attn + FFN)};
  \node[grayblk, right=0.6cm of l2,minimum width=0.6cm] (dots) {$\cdots$};
  \node[frzblk, right=0.6cm of dots,minimum width=2.0cm] (lL) {Layer $L$\\(Attn + FFN)};

  \node[trnblk,minimum width=1.6cm,above=0.6cm of emb] (p0) {$P^{(0)}$};
  \node[trnblk,minimum width=1.6cm,above=0.6cm of l1] (p1) {$P^{(1)}$};
  \node[trnblk,minimum width=1.6cm,above=0.6cm of l2] (p2) {$P^{(2)}$};
  \node[trnblk,minimum width=1.6cm,above=0.6cm of lL] (pL) {$P^{(L)}$};

  \foreach \a/\b in {emb/l1, l1/l2, l2/dots, dots/lL}
    \draw[arr] (\a)--(\b);
  \foreach \p/\l in {p0/emb, p1/l1, p2/l2, pL/lL}
    \draw[arr,green!70!black] (\p)--(\l);

  \node[green!60!black,font=\tiny,align=center] at ($(p1.east)+(0.35,0)$) {deep\\prefix};
\end{tikzpicture}
\caption{P-Tuning v2: independent trainable prefix tokens injected at every transformer layer.}
\end{figure}

\subsection{State Transition Table}
\begin{table}[h]\centering
\begin{tabular}{lll}\toprule
\textbf{Property} & \textbf{Before (Prompt Tuning)} & \textbf{After (P-Tuning v2)}\\\midrule
Prefix layers & Embedding only & All $L$ layers\\
Trainable params & $k \times d$ & $L \times l \times d$\\
NLU performance & Degrades at small scale & Competitive at all scales\\
Conditioning depth & Shallow (first layer) & Deep (every layer)\\
Task-specific storage & Single $P$ matrix & $L$ prefix matrices\\\bottomrule
\end{tabular}
\caption{State properties comparing Prompt Tuning and P-Tuning v2.}
\end{table}

